\section{Methods}

\subsection{Patient Cohort and MRI Acquisition}

Retrospective multi-b DWI data were obtained from the cohort previously described by \citet{langkilde2018evaluation}, acquired under an institutional review board--approved protocol.
For the 56 included patients, age ranged from 43 to 77 years (mean 62 years) and the mean serum prostate-specific antigen (PSA) level was 8.8~ng/mL (range, 0.37--46~ng/mL).

All patients underwent a standard multiparametric MR examination with an additional diffusion-weighted sequence on a GE 3~T Discovery MR750 scanner (GE Medical Systems, Milwaukee, WI).
An endorectal coil (Medrad, Pittsburgh, PA) was used with pelvic phased-array coils for signal reception, providing high signal-to-noise ratio at the cost of a non-routine setup.
The added DWI sequence used a matrix size of $64 \times 64$, field of view $280 \times 280\;\text{mm}^2$, slice thickness 5~mm without gap, TR~4000~ms, and TE~100~ms; the relatively long echo time was required to accommodate the high diffusion weighting with gradient times $\delta = 37$~ms and $\Delta = 45$~ms.
Fifteen b-values were acquired in 250~s/mm$^2$ increments from $b = 0$ to $3500\;\text{s/mm}^2$ along each of three orthogonal diffusion-encoding directions, combined into a trace decay (geometric mean across directions) per b-value; total scan time was 3~min.

On the multiparametric images an expert radiologist delineated, for each patient, a focal suspicious lesion (tumor) where present together with reference normal tissue in the peripheral and transition zones. For each ROI the trace signal was averaged over the segmented voxels at every b-value and paired with the patient's biopsy Gleason score.
A total of 149 ROIs were analyzed---40 tumor and 109 normal---spanning the peripheral zone (PZ; $n = 81$: 27 tumor, 54 normal) and transition zone (TZ; $n = 68$: 13 tumor, 55 normal).
Of the 40 tumor ROIs, 29 were biopsy-confirmed adenocarcinoma with a Gleason Grade Group (GGG~1, $n = 8$; GGG~2, $n = 12$; GGG~3, $n = 5$; GGG~4, $n = 2$; GGG~5, $n = 2$). The remainder were benign (7), high-grade prostatic intraepithelial neoplasia (1), or unbiopsied (3), so detection distinguishes radiologically suspicious lesions from normal tissue rather than confirmed cancer alone.

\subsection{Spectral Estimation}

For each ROI, the trace signal decay was decomposed into eight diffusivity components ($D \in \{0.25, 0.5, 0.75, 1.0, 1.5, 2.0, 3.0, 20.0\}\;\umsmm$) using both MAP and fully Bayesian estimation as described in the Theory section.

MAP spectra (Eq.~\ref{eq:map}) were computed by constrained least squares in SciPy \citep{virtanen2020scipy}.
The ridge weight $\lambda$ materially shapes the recovered spectrum: too small a value lets the ill-conditioned fit chase noise, too large a value pulls mass out of the well-determined restricted and free-water bins and spreads it across the correlated intermediate bins. We set $\lambda = 10^{-3}$ by maximizing recovery of simulated ground-truth spectra.

The posterior (Eqs.~\ref{eq:prior_R}--\ref{eq:likelihood}) was sampled independently for each of the 149 ROIs using NUTS via PyMC \citep{abrilpla2023pymc} under the half-normal spectrum prior (scale $\sigma_R = \sqrt{10} \approx 3.16$, precision $1/\sigma_R^2 = 0.1$) and a weakly informative half-Cauchy noise prior (scale $0.01$), with 2000 post-tuning draws, 200 tuning steps, 4 chains, and target acceptance rate 0.95, raised above the PyMC default to limit divergences on the correlated posterior. Sampling took 15--40~s per ROI on an 8-core Apple M1~Pro, against a sub-second MAP solve.

The noise level $\sigma$ was inferred jointly with the spectrum for each ROI rather than fixed from the acquisition; the resulting per-ROI SNR ($1/\hat{\sigma}$) had a cohort median of 303 (interquartile range 176--478).
Convergence was assessed with the rank-normalized split-$\hat{R}$ statistic \citep{vehtari2021rank} via ArviZ \citep{kumar2019arviz}: the median worst-parameter $\hat{R}$ was $1.003$, and all but two ROIs had every parameter below $1.05$; both exceptions ($1.05$, $1.06$) fell on the collinear $D = 1.5$--$3.0\;\umsmm$ bins, the least identified components. The posterior means of the spectral fractions (normalized to sum to one) were used as NUTS point estimates, and the posterior standard deviations give the per-component coefficient of variation (CV $=$ posterior std / posterior mean), used throughout as a relative identifiability measure and reported alongside the absolute posterior standard deviation (Figure~\ref{fig:fisher}b). The ordering it induces is not an artifact of fraction size: the two smallest well-identified fractions ($D = 0.25$ and $20\;\umsmm$, cohort means $0.10$ and $0.06$) have the lowest CVs, while the largest intermediate fraction ($D = 2.0\;\umsmm$, mean $0.26$) does not.

Both estimators were validated on six simulated ground-truth spectra of varying shape (smooth realistic, uniform, concentrated, and bimodal), each reconstructed at SNR 50, 300, and 1000 over 100 noise realizations (Figure~\ref{fig:validation}).
Recovery on physiologically realistic spectra was equivalent for tuned MAP and NUTS and tightened with SNR for both; a concentrated $\delta$ spike was recovered only partially and smeared into its correlated neighbors by both estimators, confirming that the smearing is the identifiability limit of the b-value grid (Theory) rather than a sampler artifact.
The jointly inferred NUTS noise level tracked the truth ($\hat{\sigma}/\sigma_{\mathrm{true}} \approx 1$) across all shapes, whereas the post-hoc MAP residual under- or over-estimated $\sigma$ depending on spectrum shape, which is why $\sigma$ is reported from the joint inference.

\subsection{ADC Computation and Classification}

Apparent diffusion coefficients were computed by log-linear least-squares fitting of $\log S(b) = \log S_0 - b \cdot \text{ADC}$ over the five b-values $\leq 1000\;\text{s/mm}^2$, following PI-RADS~v2.1 \citep{turkbey2019prostate}. The b-value range materially affects the absolute ADC value \citep{maier2022choice}.

Since normal peripheral- and transition-zone tissue differs systematically in baseline diffusion, while tumor tissue is comparatively similar across zones, tumor detection was fitted and evaluated by zone with $L_2$-regularized logistic regression under leave-one-out cross-validation (LOOCV); a pooled boundary would conflate zonal and pathological variation.
Five feature representations were scored through one common pipeline---ADC alone, each individual spectral fraction alone, the restricted and free-water fractions $\{D = 0.25,\,3.0\;\umsmm\}$ jointly, the five intermediate fractions ($D = 0.5$--$2.0\;\umsmm$), and all eight fractions (MAP or NUTS posterior means)---so that single- and multi-feature models are directly comparable. Single-feature results are quoted as raw rankings; because one standardized feature enters the logistic score monotonically, its raw-rank AUC coincides with its held-out logistic AUC.

Within each fold, features were standardized using training-ROI statistics only. Because the eight fractions differ several-fold in scale, this keeps the $L_2$ penalty from shrinking coefficients by feature scale rather than discriminative value, and renders the per-bin coefficients comparable as the change in tumor log-odds per standard deviation of each fraction.
The classifier ($C = 1.0$) was fit on the $n - 1$ remaining ROIs in the zone and applied to the held-out ROI, with all fitting confined to the training fold to prevent leakage.
ADC, alone among the representations, needs no training, so it was scored twice: as a raw ranking of ROIs by ADC value, its conventional clinical use, and through the same leave-one-out logistic pipeline as the spectral features. Only the latter is like-for-like, subjecting ADC and the spectral models to identical weight estimation and held-out scoring; the raw rank is quoted as the stricter reference (Table~\ref{tab:auc}).
Performance was the area under the receiver-operating-characteristic curve (AUC), with $95\%$ confidence intervals from patient-level percentile bootstrap (2000 resamples) \citep{efron1993bootstrap}. Sensitivity analyses at $C = 0.1$ and $C = 10$ confirmed robustness to regularization strength.
